% !TeX root = ../../portafoliofinal.tex
% !TeX encoding = utf8
% ====================================================================
%  CAPÍTULO — AUTOEVALUACIÓN DEL EQUIPO
%
%  Reflexión sobre el proceso de aprendizaje, dificultades, retos
%  superados y adquisición de competencias de la asignatura.
% ====================================================================

\chapter{Autoevaluación del equipo}

\section{Introducción}

La elaboración del portafolio grupal nos ha ofrecido la oportunidad de
reflexionar no solo sobre los contenidos teóricos de la asignatura, sino
sobre el proceso de aprendizaje en sí mismo: cómo hemos trabajado como
equipo, qué dificultades hemos encontrado, qué hemos aprendido de ellas
y en qué medida la experiencia ha contribuido a desarrollar las competencias
que la asignatura propone.

% ================================================================
\section{Lo que hemos aprendido}
% ================================================================

\subsection{Aprendizaje sobre Dirección de Recursos Humanos}

Antes de comenzar las prácticas, la mayoría del equipo tenía una visión
compartimentada de los recursos humanos: el departamento de RRHH como
una unidad que «contrata, forma y paga». El análisis de CaixaBank nos ha
revelado la dimensión estratégica y sistémica de esta función.

Algunos de los aprendizajes más significativos han sido:

\begin{itemize}
  \item \textbf{La interdependencia entre prácticas de RRHH.}
        Hemos comprendido que no existe un proceso de reclutamiento exitoso
        sin un análisis de puestos riguroso previo, ni una evaluación del
        rendimiento justa sin una retribución alineada con sus resultados.
        La elaboración del informe global fue el ejercicio que más nos hizo
        visible esta interrelación.

  \item \textbf{La brecha entre teoría y práctica empresarial.}
        Conceptos como el \textit{outplacement}, la evaluación 360° o el
        \textit{reskilling} son herramientas teóricamente robustas que, en
        la práctica de una empresa real, se aplican de forma parcial,
        adaptada o directamente no se aplican. El caso del \textit{outplacement}
        en CaixaBank (disponible solo puntualmente, no como política estable)
        fue revelador.

  \item \textbf{El peso de la negociación sindical.}
        El ERE de 2021 nos enseñó que las decisiones de RRHH en grandes
        organizaciones nunca son puramente técnicas: el marco sindical,
        los acuerdos colectivos y la presión social condicionan de forma
        determinante lo que es posible o deseable hacer.

  \item \textbf{La transformación digital como eje transversal.}
        La digitalización no es solo un tema tecnológico; es un reto de
        gestión del talento. Hemos aprendido que sustituir perfiles
        tradicionales por perfiles digitales requiere planificación,
        formación, gestión de la salida y comunicación interna de manera
        simultánea y coordinada.
\end{itemize}

\subsection{Aprendizaje sobre CaixaBank}

Más allá de la teoría, el análisis en profundidad de CaixaBank nos ha
proporcionado un conocimiento concreto y aplicado del sector bancario español.
Hemos entendido cómo opera la mayor entidad bancaria española en materia de
capital humano, qué decisiones ha tomado en momentos críticos (fusión con Bankia)
y qué retos afronta en los próximos años.

% ================================================================
\section{Dificultades encontradas}
% ================================================================

\subsection{Acceso a información interna}

La principal dificultad del proyecto ha sido la limitación en el acceso a
datos internos de CaixaBank. Las personas entrevistadas —por política de
confidencialidad y por el nivel de responsabilidad de su cargo— no pudieron
compartir información cuantitativa sobre presupuestos de formación, costes
de reclutamiento o KPIs concretos del departamento de Personas.

Esta limitación nos obligó a desarrollar la capacidad de analizar con datos
incompletos: triangular fuentes públicas (Informe Anual, Plan Estratégico,
informe de sostenibilidad) con la información cualitativa obtenida en las
entrevistas y con el marco teórico de la asignatura.

\subsection{Coordinación del trabajo en equipo}

Coordinar un equipo de cinco personas con diferentes horarios, asignaturas
compartidas y estilos de trabajo ha requerido una gestión activa del proceso.
Hemos tenido que tomar decisiones sobre cómo dividir el trabajo sin perder
la coherencia del conjunto, y cómo integrar contribuciones con distintos
niveles de profundidad y extensión.

La distribución del trabajo en bloques temáticos individuales facilitó la
división de tareas, pero dificultó la integración final, que requirió
revisiones conjuntas para asegurar que los capítulos se complementaban
y no se contradecían entre sí.

\subsection{Uso del lenguaje técnico y herramientas de redacción}

La redacción académica en un documento LaTeX ha representado un reto técnico
para varios miembros del equipo. Hemos invertido tiempo en aprender a
estructurar tablas, referencias cruzadas y el sistema de citas, lo que inicialmente
ralentizó el proceso pero a largo plazo ha supuesto una herramienta de valor.

% ================================================================
\section{Retos superados}
% ================================================================

\begin{itemize}
  \item \textbf{Conseguir tres entrevistas con empleados de CaixaBank.}
        No fue sencillo encontrar personas dispuestas a dedicar tiempo y
        compartir información sobre procesos internos. Hemos aprendido a
        redactar solicitudes de entrevista profesionales y a adaptar nuestro
        guion de preguntas cuando el entrevistado tenía información parcial
        sobre algunos temas.

  \item \textbf{Sintetizar grandes volúmenes de información.}
        CaixaBank publica una cantidad ingente de información corporativa.
        Aprender a identificar lo relevante, descartar lo accesorio y
        estructurar una narrativa coherente ha sido uno de los ejercicios
        de síntesis más exigentes de nuestra trayectoria académica hasta
        el momento.

  \item \textbf{Conectar teoría y práctica de forma crítica.}
        Ha sido fácil caer en la tentación de describir las prácticas de
        CaixaBank como ejemplos perfectos de lo estudiado en clase. El
        verdadero aprendizaje ha venido de la mirada crítica: identificar
        dónde CaixaBank se queda corta respecto a sus propios estándares
        declarados y proponer mejoras fundamentadas.

  \item \textbf{Elaborar el informe global.}
        Fue el apartado que más tiempo requirió, precisamente porque no
        se trataba de repetir lo ya dicho sino de ver el conjunto desde
        un nivel de abstracción mayor. Las primeras versiones eran
        recopilaciones de recomendaciones anteriores. La versión final,
        con recomendaciones sistémicas que interrelacionan múltiples
        prácticas, fue posible solo después de releer todo el trabajo y
        discutir en equipo qué hilos conectores eran realmente relevantes.
\end{itemize}

% ================================================================
\section{Valoración de la aplicación práctica de la asignatura}
% ================================================================

La metodología de trabajo práctico sobre una empresa real es, a nuestro
juicio, la parte más valiosa de la asignatura. Permite que los conceptos
teóricos tengan una referencia concreta, que las limitaciones de la teoría
sean visibles y que el aprendizaje tenga un anclaje en la realidad que
los exámenes o ejercicios teóricos no proporcionan.

En particular, valoramos positivamente:

\begin{itemize}
  \item La posibilidad de contactar con profesionales reales de la empresa,
        lo que añade autenticidad al análisis y desarrolla habilidades de
        comunicación profesional.
  \item La estructura acumulativa de las tres entregas, que permite
        profundizar progresivamente en la empresa y construir un conocimiento
        integrado a lo largo del semestre.
  \item La exigencia del informe global, que obliga a salir del análisis
        parcial y desarrollar una visión sistémica de la gestión del talento.
\end{itemize}

Como aspecto mejorable, señalamos que sería útil disponer de un mayor
acompañamiento en la fase de elaboración del informe global, dado que es
el apartado más complejo conceptualmente y en el que la guía de los criterios
de evaluación es menos explícita.

% ================================================================
\section{Adquisición de competencias de la asignatura}
% ================================================================

A continuación valoramos de forma reflexiva en qué medida este trabajo
ha contribuido al desarrollo de las competencias específicas de la asignatura:

\begin{table}[H]
  \centering
  \small
  \caption{Valoración de la adquisición de competencias transversales}
  \begin{tabularx}{\textwidth}{p{0.40\textwidth} X}
    \toprule
    \textbf{Competencia}                         & \textbf{Valoración y evidencia} \\
    \midrule
    \textbf{Trabajo en equipo}
    & Alta. La coordinación de cinco personas durante un semestre,
      con roles diferenciados pero responsabilidad compartida sobre
      el resultado final, ha sido el ejercicio de trabajo cooperativo
      más exigente de nuestra trayectoria hasta la fecha. \\
    \addlinespace
    \textbf{Liderazgo y gestión de roles}
    & Media-alta. El equipo adoptó naturalmente una distribución de
      roles de coordinación, redacción y revisión sin una estructura
      formal. Identificar y respetar las fortalezas individuales fue
      clave para avanzar sin conflictos. \\
    \addlinespace
    \textbf{Gestión del tiempo}
    & Media. Las tres entregas con plazos definidos nos han obligado a
      planificar el trabajo con antelación. Sin embargo, reconocemos que
      en las entregas iniciales subestimamos el tiempo necesario para la
      integración final de los bloques individuales. \\
    \addlinespace
    \textbf{Capacidad de síntesis}
    & Alta. Resumir la política de RRHH de una empresa con 45.000 empleados
      en un documento académico estructurado ha exigido un ejercicio
      constante de síntesis, selección y priorización de la información. \\
    \addlinespace
    \textbf{Adaptación de contenido teórico a la realidad empresarial}
    & Alta. Esta competencia ha sido, a nuestro juicio, la más desarrollada.
      La comparación sistemática entre el marco teórico y la práctica real
      de CaixaBank —con sus contradicciones, limitaciones y logros— ha
      enriquecido notablemente nuestra comprensión de la disciplina. \\
    \bottomrule
  \end{tabularx}
\end{table}

\medskip

En resumen, la experiencia de elaborar este portafolio ha superado nuestras
expectativas iniciales. Comenzamos el curso con la idea de que «recursos humanos»
era un campo relativamente acotado de la gestión empresarial; lo concluimos con
la certeza de que es, en realidad, uno de los ámbitos más complejos, transversales
y estratégicamente relevantes de cualquier organización.

\endinput
